% forma de utilizar no texto:
% \ac{ages} << exemplo

\begin{center}
  \uppercase{\bfseries lista de siglas}\\[3em]
\end{center}

\begin{acronym}[TECNOPUC]
  % Só vai aparecer no texto se for utilizado
  % \acro{chave}[SIGLA]{Descrição completa}
  \acro{abnt}   [ABNT]   {Associação Brasileira de Normas Técnicas}
  \acro{adm}    [ADM]    {Administrador}
  \acro{ages}   [AGES]   {Agência Experimental de Engenharia de Software}
  \acro{api}    [API]    {Application Programming Interface}
  \acro{aws}    [AWS]    {Amazon Web Services}
  \acro{crud}   [CRUD]   {Create, Read, Update, Delete}
  \acro{css}    [CSS]    {Cascading Style Sheets}
  \acro{dto}    [DTO]    {Data Transfer Object}
  \acro{ens}    [ENS]    {Escola de Raquetes ENSportive}
  \acro{html}   [HTML]   {HyperText Markup Language}
  \acro{ia}     [IA]     {Inteligência Artificial}
  \acro{jdbc}   [JDBC]   {Java Database Connectivity}
  \acro{mr}     [MR]     {Merge Request}
  \acro{mvp}    [MVP]    {Minimum Viable Product}
  \acro{ong}    [ONG]    {Organização Não Governamental}
  \acro{pcd}    [PCD]    {Pessoa com Deficiência}
  \acro{poc}    [POC]    {Proof of Concept}
  \acro{pucrs}  [PUCRS]  {Pontifícia Universidade Católica do Rio Grande do Sul}
  \acro{rest}   [REST]   {Representational State Transfer}
  \acro{sql}    [SQL]    {Structured Query Language}
  \acro{uml}    [UML]    {Unified Modeling Language}
\end{acronym}
